\documentclass[11pt]{article}
\usepackage[T1]{fontenc}
\usepackage[margin=1in]{geometry}
\usepackage{microtype}
\usepackage{amsmath,amssymb}
\usepackage{booktabs}
\usepackage{siunitx}
\usepackage{hyperref}
\usepackage{graphicx}
\usepackage{epstopdf}
\usepackage{caption}
\usepackage{xcolor}
\usepackage{array}
\usepackage{longtable}
\usepackage{enumitem}

\sisetup{
  separate-uncertainty = true,
  per-mode = symbol,
}

\hypersetup{
  colorlinks=true,
  linkcolor=blue!60!black,
  citecolor=blue!60!black,
  urlcolor=blue!60!black,
}

\title{FELsim Beamline Optimisation:\\
  Technical Compendium}
\author{Eremey Valetov}
\date{11 March 2026}

\begin{document}
\maketitle

\begin{abstract}
We present a comprehensive summary of the FELsim transport line
optimisation programme for the UH MkV free electron laser.  The programme
encompasses 28~completed studies across five categories: parameter
sensitivity (S1--S9), cross-code validation (C1--C8), optimizer
improvements (O1--O5), physics model upgrades (P1--P3), and infrastructure
(I1--I8).  We validated the 11-stage Nelder-Mead optimisation against
COSY~INFINITY (differential algebra) and RF-Track (particle tracking),
achieving cross-code MSE agreement to $10^{-6}$--$10^{-9}$.  A two-phase
NM$\to$CMA-ES polishing strategy improves Stage~11 precision by up to
$10^{10}\times$.  Chromatic physics (per-particle momentum-dependent
matrices for quads, dipoles, and edge kicks) and aperture loss tracking
were integrated into the production optimizer, raising simulated
transmission from 64.5\% to 98.1\%.  The feasibility boundary is dominated
by emittance; energy spread and chirp have negligible impact on
transverse matching at the baseline $\varepsilon_n = 8\;
\pi\cdot\text{mm}\cdot\text{mrad}$.  The transport line cannot compress
bunches below $\sim$\SI{0.45}{ps} (limited by $R_{56} \times
\sigma_\delta$); compression must occur upstream.
\end{abstract}

\tableofcontents
\newpage


%=======================================================================
\section{Introduction}
%=======================================================================

The UH Mark~V FEL transport line consists of 118~beamline elements
(23~quadrupoles, 4~dipole chicane sections, drifts, correctors, and
diagnostics) that transport a \SI{40}{MeV} electron beam from the linac
exit to the undulator entrance.  The beam must be matched to the
undulator Twiss targets specified by Weinberg, Fisher \&
Li~\cite{weinberg2025}:
\begin{equation}
  \beta_x = \SI{1.400}{m},\quad
  \alpha_x = 0.470,\quad
  \beta_y = \SI{0.242}{m},\quad
  \alpha_y = 0.
  \label{eq:targets}
\end{equation}

FELsim decomposes this 23-variable matching problem into an 11-stage
sequential optimisation, where each stage controls a physically motivated
group of quadrupoles (singlets, doublets, triplets, chromaticity
sections).  The objective function is the weighted mean squared error
(MSE) between computed and target Twiss parameters.  Results are reported
as the root mean square (RMS) $= \sqrt{\text{MSE}}$, which has the same
units as the Twiss residuals.

This compendium summarises all completed studies, organised by topic
rather than chronology, and provides cross-references to the detailed
scripts, reports, and result files.


%=======================================================================
\section{Simulation Framework}
\label{sec:framework}
%=======================================================================

\subsection{FELsim (First-Order Transfer Matrices)}

FELsim's native transport model uses $6\times6$ linear transfer matrices
(\texttt{beamline.py}).  Each element type has an analytical matrix:
drifts, thick quadrupoles with $\cos$/$\cosh$ focusing, sector-bend
dipole bodies, and dipole wedges with Enge-function fringe corrections.
The model evaluates in $\sim$\SI{1}{ms} per beamline, enabling
$\mathcal{O}(10^4)$ optimizer evaluations per study.

\paragraph{Chromatic extensions (P1--P3).}
All elements support an opt-in \texttt{chromatic=True} mode that computes
per-particle momentum-dependent matrices:
\begin{itemize}[nosep]
  \item \textbf{Quadrupoles (C8):} $k_\text{eff} = k_0 \times
    \beta\gamma_0 / \beta\gamma(\delta)$, matching RF-Track's intrinsic
    momentum scaling to 0.00\% across all Twiss parameters.
  \item \textbf{Dipole body (P3):} Per-particle $\rho_p = \rho_0 \times
    \beta\gamma_p / \beta\gamma_0$, $\theta_p = L/\rho_p$.
  \item \textbf{Dipole wedge (P2):} Momentum-dependent $R$, $\phi$, and
    edge kicks ($h_p = h_0 \times P_0/P$).
\end{itemize}
All chromatic matrices are symplectic ($\det M = 1.0$ to machine
precision) and vectorised over particles without Python loops.

\paragraph{Aperture loss tracking (P1).}
Physical apertures are defined per element type: QPF/QPD
$\pm$\SI{13.5}{mm} (bore), DPH/DPW $\pm$pole\_gap/2.  The optimizer
integration adds a smooth $W \times (1-T)^2$ penalty to the MSE
objective when aperture tracking is enabled.


\subsection{COSY INFINITY (Differential Algebra)}

COSY~INFINITY computes transfer maps to arbitrary order using
differential algebra (DA).  The Python adapter (\texttt{cosyAdapter.py})
generates FOX code, executes the COSY binary, and parses output.
Quadrupoles are thick \texttt{MQ} elements; dipoles are thick
\texttt{DIL} with integrated edge angles and optional Enge fringe fields
(\texttt{FC}/\texttt{FD}) or measured fieldmaps (\texttt{MGE}).  The
built-in FIT optimizer exploits DA derivatives for gradient-based matching.


\subsection{RF-Track (Particle Tracking)}

RF-Track (v2.5.5, CERN) performs element-by-element particle tracking
with support for space charge.  The adapter (\texttt{rftrackAdapter.py})
includes an analytical sector-bend workaround for the \texttt{SBend}
P/$\delta$ bug (C5), thin-lens quadrupole edge kicks for DPW elements,
and prefix caching for fast Stage~11 re-tracking.


\subsection{Multi-Code Orchestrator (I7)}

The \texttt{MultiCodeSimulator} chains multiple adapters on contiguous
beamline sections, using FELsim coordinates as the interchange format.
Validated configurations include FELsim$\to$RF-Track, FELsim$\to$COSY,
COSY$\to$RF-Track, and three-way FELsim$\to$COSY$\to$RF-Track hybrids.


%=======================================================================
\section{Transverse Beam Matching}
\label{sec:transverse}
%=======================================================================

\subsection{Baseline Optimisation (S1, S3)}

The 11-stage Nelder-Mead (NM) optimiser achieves MSE $< 10^{-5}$ at the
baseline operating point ($\varepsilon_n = 8\;\pi\cdot\text{mm}\cdot\text{mrad}$,
$\sigma_E = 0.5\%$, $h = 5 \times 10^9\;\text{s}^{-1}$, 0.5~ps bunch).
The 2~ps (S1) and 0.5~ps (S3) optimisations produce identical quadrupole
currents, confirming transverse--longitudinal decoupling at first order
(S9).

\subsection{Cross-Code Validation (W4, C1, R2)}

Table~\ref{tab:xcode} summarises the cross-code comparison at the
baseline operating point.

\begin{table}[ht]
\centering
\caption{Cross-code optimisation at $\varepsilon_n = 8\;
  \pi\cdot\text{mm}\cdot\text{mrad}$.  Targets:
  $\beta_x = \SI{1.400}{m}$, $\alpha_x = 0.470$,
  $\beta_y = \SI{0.242}{m}$, $\alpha_y = 0$.}
\label{tab:xcode}
\begin{tabular}{@{} l l S[scientific-notation=true] @{}}
\toprule
Code & Model & {MSE} \\
\midrule
FELsim NM       & Transfer matrix       & 1.2e-06 \\
COSY FR\,0      & DA map (no fringe)    & 2.3e-07 \\
COSY FR\,1      & DA map (1st fringe)   & 3.7e-09 \\
RF-Track opt    & Particle tracking     & 7.0e-03 \\
\bottomrule
\end{tabular}
\end{table}

FELsim and COSY achieve MSEs of $10^{-6}$--$10^{-9}$, matching all four
Twiss targets to within $10^{-3}$.  COSY's gradient-based DA optimizer
converges $\sim 10^3\times$ more precisely.  RF-Track's higher MSE
($7.0 \times 10^{-3}$) reflects a systematic $\beta_y$ deficit from the
thin-lens edge-kick model (missing triangle-rule fringe correction
$\phi$), not optimizer limitations.

\paragraph{Bug fixes (R2).}
Achieving cross-code agreement required identifying and correcting
9~implementation errors: FELsim DPW edge-kick sign inversion, five
RF-Track adapter issues (SBend P/$\delta$ confusion, DPW modelling,
edge-kick sign, R-matrix units, coord5 conversion), and three COSY
adapter issues (initial $\beta_0$, FIT objective reset, warm-start for
fringe fields).  See the R2 report for details.


\subsection{Algorithm Comparison}
\label{sec:algorithms}

\paragraph{Nelder-Mead (NM).}
The baseline optimizer.  Multi-start with 5~random Stage~11 initial
conditions resolves local minimum traps at extreme emittances
($\varepsilon_n = 14$--16).

\paragraph{Glyfada (W6).}
The glyfada evolutionary optimizer (CMA-ES, NSGA-II) via the distributed
DH evaluator protocol.  At 600~evaluations per run, glyfada fails at all
three test emittances ($\varepsilon_n = 5, 8, 14$); NM outperforms by
3--6~orders of magnitude in MSE.  Root cause: uniform random
initialisation over wide bounds wastes evaluations in infeasible regions.

\paragraph{CMA-ES polishing (W7, O5).}
pycma CMA-ES warm-started from the NM solution dramatically outperforms
NM alone:

\begin{table}[ht]
\centering
\caption{Two-phase NM$\to$CMA-ES vs NM-only for Stage~11 matching.}
\label{tab:cmaes}
\begin{tabular}{@{} l S[scientific-notation=true] S[scientific-notation=true] l @{}}
\toprule
$\varepsilon_n$ & {NM MSE} & {CMA-ES MSE} & Improvement \\
\midrule
5   & 8.2e-06 & 4.5e-16 & $\sim 10^{10}\times$ \\
8   & 6.1e-06 & 1.2e-15 & $\sim 10^{9}\times$ \\
14  & 8.9e-03 & 6.3e-03 & $1.4\times$ \\
\bottomrule
\end{tabular}
\end{table}

The two-phase strategy was integrated into the production optimizer
(O5): after NM converges, CMA-ES refines from the NM solution with
$\sigma = 0.1$, popsize$= 20$, maxfevals $= 3000$.  The composite
5-goal objective (4~Twiss + dispersion) has a dispersion floor of
$\sim 2.5 \times 10^{-6}$ that Stage~11 quads cannot reduce.

\paragraph{Stage~11 unimodality (S8).}
A multi-start robustness study (5~emittance points $\times$
10~random starts) found that \emph{all} random starts converge to
identical solutions at machine precision.  The Stage~11 landscape has a
single global basin; multi-start cannot improve it.  The
seed-dependence observed at extreme emittances (S7) originates entirely
from stages~1--10 (random beam generation affecting upstream matching).

\paragraph{Warm-starting (O1).}
Sequential warm-starting from neighbouring scan points is
counterproductive: the $\varepsilon_n = 1$ basin (Marginal, RMS
$\approx 0.108$) traps the optimizer, propagating poor solutions through
$\varepsilon_n = 2$--13.  Cold start independently finds better basins at
each point.  Warm-start wins only 4/20~emittance points.


%=======================================================================
\section{Parameter Sensitivity}
\label{sec:sensitivity}
%=======================================================================

\subsection{1D Scans (S4)}

Three 1D sweeps at the 0.5~ps operating point (500~particles, seed~42):
\begin{itemize}[nosep]
  \item \textbf{Energy spread} ($\sigma_E = 0.1$--5\%): all Excellent
    ($\text{MSE} < 10^{-3}$) up to $\sigma_E \approx 3\%$.
  \item \textbf{Chirp} ($h = 0$--$40 \times 10^9$~s$^{-1}$): chirp has
    negligible effect on transverse matching (all Excellent).
  \item \textbf{Emittance} ($\varepsilon_n = 1$--20): the dominant
    sensitivity parameter.  Fails at $\varepsilon_n \leq 2$ (beam too
    small) and $\varepsilon_n \geq 17$ (requires raised bounds).
\end{itemize}

\subsection{2D Scans (S5)}

Three $10\times10$ grids (100~points each, 500~particles):
\begin{itemize}[nosep]
  \item \textbf{S5a} ($\sigma_E \times h$) at $\varepsilon_n = 8$:
    0~failures, 92~Excellent.  Energy spread and chirp do not compromise
    matching.
  \item \textbf{S5b} ($\sigma_E \times \varepsilon_n$) at $h = 5 \times
    10^9$: 18~failures at extreme emittances.
  \item \textbf{S5c} ($h \times \varepsilon_n$) at $\sigma_E = 0.5\%$:
    20~failures at extreme emittances.  Chirp does not rescue
    high-emittance failure.
\end{itemize}
The feasibility boundary is dominated by emittance.

\subsection{Bunch Length Independence (S6, S9)}

The 2~ps and 0.5~ps operating points produce identical quad currents.
Root cause: transverse Twiss depends on $\varepsilon$ and $\sigma_\delta$,
not $\sigma_z$; bunch length enters only column~5 (time of flight), which
\texttt{cal\_twiss()} ignores.  A 15-point bunch-length sweep
(\SIrange{0.1}{2.0}{ps}) confirms no degradation at any bunch length.

\subsection{Verification (S7)}

Re-running 7~emittance points and 3~energy-spread points at
$N = 500, 1000, 2000$ particles reveals:
\begin{itemize}[nosep]
  \item Energy spread results are fully consistent across particle counts.
  \item Emittance results at extremes ($\varepsilon_n \neq 8$) are
    \emph{not} robust --- different particle samples lead to different NM
    basins in stages~1--10.  Publication-quality claims at extreme
    emittances require multi-start + multi-seed statistics.
\end{itemize}


%=======================================================================
\section{Longitudinal Dynamics}
\label{sec:longitudinal}
%=======================================================================

\subsection{Transfer Map Analysis (W9)}

COSY~INFINITY 6D maps yield:
$R_{56} = \SI{+27.09}{mm}$ (elongates unchirped beams),
$T_{566} = 0$ (no second-order path-length dependence).
The transport line is intrinsically non-compressing.

\subsection{Bunch Compression Feasibility (W12)}

Can the transport line compress 2~ps $\to$ 0.5~ps?
\begin{itemize}[nosep]
  \item \textbf{Analytical compression}: floor $\approx$~\SI{0.45}{ps}
    from $R_{56} \times \sigma_\delta$ at $\sigma_E = 0.5\%$.
  \item \textbf{Chirp}: $C = 4$ chirp gives $\sim$\SI{0.67}{ps}, still
    above target.
  \item \textbf{Extended bounds} (15~A, 5~restarts): improves transverse
    MSE; $\sigma_z$ unchanged.
  \item \textbf{Conclusion}: the transport line is not a compressor.
    Compression must occur upstream (velocity bunching or dedicated
    compressor).
\end{itemize}

\subsection{$\sigma_z$ Blowup Investigation (I6)}

The 60--100$\times$ longitudinal blowup originally reported was
\emph{not} reproduced with properly matched optics (MSE $\approx 2.3
\times 10^{-7}$).  With optimised currents: ORDER~1 gives $\sigma_z =
\SI{2.264}{ps}$ ($1.13\times$), ORDER~3 gives $\sigma_z =
\SI{4.158}{ps}$ ($2.08\times$).  The higher-order growth is correct
physics from nonlinear path-length terms.


%=======================================================================
\section{Chromatic Optimisation}
\label{sec:chromatic}
%=======================================================================

Linear-optics matching ignores the momentum dependence of focusing
strengths.  At $\sigma_p = 0.5\%$, linear FELsim shows $\beta_y$ errors
up to 1466\% compared with RF-Track's intrinsic chromatic tracking (C8).
After enabling chromatic mode for all elements (quads, dipole bodies,
dipole wedges), a three-pass coordinate-descent re-optimisation with
aperture losses yields:
\begin{itemize}[nosep]
  \item MSE $= 4.0 \times 10^{-4}$
    ($\beta_x = 1.40$, $\alpha_x = 0.44$, $\beta_y = 0.23$,
    $\alpha_y = 0.03$)
  \item Transmission $= 98.1\%$ (up from 64.5\% with linear-optimised
    currents)
  \item Significant current changes in the chicane region (elements
    33--43) and downstream quads (70, 87, 93--97)
\end{itemize}


%=======================================================================
\section{Beam Losses and Throughput}
\label{sec:losses}
%=======================================================================

\subsection{Loss Distribution (P1)}

FELsim aperture tracking at the baseline operating point
(2000~particles, $\sigma_p = 0.5\%$):
\begin{itemize}[nosep]
  \item Overall transmission: 64.5\% (with linear-optimised currents)
  \item Dominant loss: FC1.DPW.111 (element 100, $\pm$\SI{6.3}{mm})
    --- 395~particles (56\% of losses)
  \item Chicane exit dipoles (72--73, 89--90): 281~particles
  \item Quad losses minor: 34~particles total
\end{itemize}

\subsection{Throughput Optimisation (W11)}

Extending the Stage~11 objective to include transmission and bunch length
via weighted scalar cost:
$w_t \cdot \text{MSE}_\text{Twiss} + w_T \cdot (1-T)^2 + w_\sigma
\cdot (\sigma_t/\sigma_\text{target} - 1)^2$.
Results (RF-Track with physical apertures):
\begin{itemize}[nosep]
  \item 2~ps transport: MSE $= 0.175$, $T = 42.8\%$, $\sigma_t =
    \SI{1.83}{ps}$, $I_\text{peak} = \SI{5.6}{A}$
  \item 0.5~ps compression: MSE $= 0.071$, $T = 28.4\%$, $\sigma_t =
    \SI{1.31}{ps}$ (target 0.5), $I_\text{peak} = \SI{5.2}{A}$
\end{itemize}
Compression failed --- $\sigma_t$ target unreachable with Stage~11 quads
alone, consistent with the W12 feasibility finding.


%=======================================================================
\section{FR3+MGE Optimisation (C3)}
\label{sec:mge}
%=======================================================================

Matching the beamline with COSY~INFINITY's 3rd-order fringe fields plus
measured fieldmaps (MGE) requires simultaneous optimisation of all
23~quadrupole currents.  Sequential FIT fails because the per-stage local
optima under FR3+MGE are incompatible with the global solution.

Four CMA-ES campaigns on the Koa HPC cluster:
\begin{itemize}[nosep]
  \item \textbf{v1}: 10,020~evals, MSE$=1030$ (barely unstable,
    $\cos\mu \approx 1.03$).
  \item \textbf{v2}: 16,000~evals, sigma collapsed to $1.3 \times
    10^{-5}$; no stable solution found.
  \item \textbf{v3}: Killed after 247~evals.
  \item \textbf{v4} (current): Capped objective ensuring any stable
    solution beats any unstable one.  663~evals completed, best objective
    12,009 (unstable but steadily improving).
\end{itemize}
Root cause: the MSE landscape has an extremely narrow stability basin in
23D.  The capped objective (v4) directly addresses the pathology where
unstable solutions could outcompete stable ones with poor Twiss.


%=======================================================================
\section{Infrastructure}
\label{sec:infrastructure}
%=======================================================================

\paragraph{Element labels (I1).}
All beamline elements carry human-readable names propagated from Excel
(Label/Nomenclature columns) and JSON/YAML (\texttt{name} field) through
the entire pipeline.

\paragraph{PALS alignment (I3, I8).}
The v3 lattice format adds PALS-aligned fields (\texttt{Bn1},
\texttt{BendP}) with full backward compatibility.  The standalone
\texttt{pals2cosy} converter (v0.3.0) supports both official PALS and
FELsim v2/v3 formats.

\paragraph{Test suite (C4, I6).}
210~tests (7~skipped), covering: unit tests for transfer matrices and
Twiss computation; frozen regression tests for the full beamline;
cross-code benchmarks (FELsim vs RF-Track); adapter round-trips
(Excel/JSON/YAML); chromatic physics (30~tests); attribute guards
(AST-based typo detection).

\paragraph{Multi-code framework (I7).}
The \texttt{MultiCodeSimulator} enables hybrid simulations with per-section
code selection and configuration passthrough (space charge, apertures).
Validated for FELsim+RF-Track, FELsim+COSY, COSY+RF-Track, and three-way
hybrids.


%=======================================================================
\section{Summary of Findings}
\label{sec:summary}
%=======================================================================

\begin{enumerate}[nosep]
  \item The 11-stage sequential NM optimisation reliably matches the
    undulator Twiss targets at the baseline operating point, with
    cross-code agreement to MSE $\sim 10^{-6}$--$10^{-9}$.
  \item Two-phase NM$\to$CMA-ES polishing improves Stage~11 precision by
    up to $10^{10}\times$ and is integrated as the default post-processing
    step.
  \item The feasibility boundary is dominated by normalised emittance
    ($\varepsilon_n$).  Energy spread ($\sigma_E$) and chirp ($h$) have
    negligible impact at the baseline $\varepsilon_n = 8\;
    \pi\cdot\text{mm}\cdot\text{mrad}$.
  \item Bunch length does not affect transverse matching (confirmed
    analytically and numerically).
  \item Bunch compression below $\sim$\SI{0.45}{ps} is not achievable
    by the transport line alone ($R_{56} = +\SI{27}{mm}$, $T_{566} = 0$).
  \item Chromatic optimisation with aperture tracking raises transmission
    from 64.5\% to 98.1\% with MSE $= 4.0 \times 10^{-4}$.
  \item The Stage~11 landscape is unimodal; robustness issues at extreme
    emittances originate from upstream stages, not the final matching
    section.
  \item Sequential warm-starting is counterproductive for this beamline
    (basin trapping).
  \item FR3+MGE optimisation requires simultaneous 23-variable search
    (in progress on HPC).
\end{enumerate}


%=======================================================================
\section{Study Index}
\label{sec:index}
%=======================================================================

\small
\begin{longtable}{@{} l l l l @{}}
\toprule
ID & Title & Status & Key Result \\
\midrule
\endfirsthead
\toprule
ID & Title & Status & Key Result \\
\midrule
\endhead
\bottomrule
\endlastfoot
%
S1 & 2 ps Baseline Optimisation & Done & MSE $< 10^{-5}$ \\
S2 & 0.5 ps Emittance Conservation & Done & Joint 4-var Stage~11 \\
S3 & 0.5 ps Paper-Aligned & Done & Identical to S1 currents \\
S4 & 1D Parameter Sensitivity & Done & Emittance dominates \\
S5 & 2D Coupled Scans & Done & $\sigma_E$, $h$ negligible \\
S6 & Bunch Length Sensitivity & Done & No effect on matching \\
S7 & Verification Runs & Done & Extremes not robust \\
S8 & Multi-Start Robustness & Done & Stage~11 unimodal \\
S9 & Bunch Length Independence & Done & Linear decoupling \\
\midrule
W1 & Chirp Comparison & Done & $h$ negligible on Twiss \\
W2 & Emittance Multi-Start & Done & Dips at $\varepsilon_n = 14$--16 resolved \\
W4 & COSY Cross-Validation & Done & FR\,1 MSE $= 3.7 \times 10^{-9}$ \\
W5 & Chirp Assessment & Done & $h = 5 \times 10^9$ is low-moderate \\
W6 & Glyfada Benchmark & Done & NM wins by $10^{3}$--$10^{6}\times$ \\
W7 & CMA-ES Re-Benchmark & Done & CMA-ES wins by $10^{9}$--$10^{10}\times$ \\
W9 & COSY Longitudinal & Done & $R_{56} = +27$ mm, $T_{566} = 0$ \\
W10 & Beam Losses Study & Done & 40--50\% loss at 0.5 ps \\
W11 & Throughput Optimisation & Done & $T = 42.8\%$ at 2 ps \\
W12 & Compression Feasibility & Done & Floor $\approx$ 0.45 ps \\
\midrule
C1 & RF-Track Cross-Validation & Done & 5-way comparison \\
C3 & FR3+MGE Optimisation & In Progress & CMA-ES v4 on Koa \\
C5 & RF-Track SBend Bug & Done & Analytical workaround \\
C8 & Chromatic Quad Matrices & Done & 0.00\% vs RF-Track \\
\midrule
O1 & Warm-Starting & Done & Counterproductive \\
O3 & Glyfada Adapter & Done & \texttt{method='glyfada'} \\
O5 & CMA-ES Polishing & Done & Integrated in production \\
\midrule
P1 & Aperture Loss Tracking & Done & 64.5\% baseline $T$ \\
P2 & Chromatic DPW Edges & Done & Symplectic, vectorised \\
P3 & Chromatic Dipole Body & Done & $< 1\%$ at $\sigma_p = 0.5\%$ \\
\midrule
I1 & Element Labels & Done & End-to-end propagation \\
I3 & PALS2COSY Converter & Done & v0.3.0 (55 tests) \\
I7 & Multi-Code Framework & Done & 3-way hybrid validated \\
I8 & v3 Lattice Format & Done & PALS-aligned fields \\
\end{longtable}
\normalsize


\begin{thebibliography}{9}
\bibitem{weinberg2025}
  A.~Weinberg, A.~Fisher, and R.~Li,
  ``Design of a compact FEL oscillator for the University of Hawai`i,''
  arXiv:2510.14061v1, 2025.
\end{thebibliography}

\end{document}
